\section{The critical block tree}\label{block:section}

We now pass from the word estimates to an exact representation of the
finite map law. A \emph{corner} of a nonempty map is a sector between
successive darts at a vertex. Thus a map with $m$ edges has $2m$
corners, including two at a vertex carrying a single loop.

\subsection{Decomposition at corners}

The decomposition of a rooted map into a root block and corner insertions
goes back to Tutte~\cite[Section~6]{Tutte1963}. Addario-Berry makes its
even-offspring tree representation explicit for uniform
maps~\cite[Section~2 and Proposition~3.1]{AddarioBerry2019}; the
block-weighted form belongs to Stufler's enriched-tree
framework~\cite[preprint Section~6.1.5]{Stufler2020}.
We prove the rooted version here to keep the one-link and one-loop
conventions, trivial inserts, and exact finite weights explicit.

Our nontrivial \emph{blocks} are the rooted one-loop map and the rooted
nonempty connected loopless maps with at least two vertices and no
separating vertex. The one-link map is included: deleting either of its
vertices leaves a single vertex. Write $b_m$ for the total weight
$a(B)$ of the blocks with $m$ edges. In particular,
\begin{equation}\label{block:one-edge}
 b_1=2(1+t)>0.
\end{equation}

\begin{lemma}[Rooted corner substitution]\label{block:decomposition}
Every nontrivial rooted planar map is uniquely obtained from a rooted
block by inserting a rooted map, possibly trivial, in each of its
corners. A nonempty insert is attached at the origin of its root dart.
Moreover, $a$ is multiplicative when two maps are joined at one vertex.
Consequently, with
\[
 T(z)=\sum_{n\ge0}a_nz^n,\qquad
 \Phi(y)=1+\sum_{m\ge1}b_my^m,
\]
there is an identity of formal power series
\begin{equation}\label{block:equation}
 T(z)=\Phi\big(zT(z)^2\big).
\end{equation}
\end{lemma}

\begin{proof}
Suppose first that the root is not a loop. Take the maximal loopless
block containing the root edge. It is unique: the union of two
connected subgraphs without a separating vertex that share an edge
again has no separating vertex. An exterior path joining two different
vertices of this block, with interior disjoint from it, could be added
without producing a separating vertex. Thus every component outside the
block attaches at one block vertex.

Each face boundary of a loopless block visits a given vertex at most
once. Indeed, if two distinct sectors at a vertex belonged to the same
face, a simple curve through that face joining the sectors and closed
at the vertex would separate incident edges on its two sides.
Connectivity after deleting the vertex rules this out. The assertion
also holds for a one-link block. Planarity therefore places each
outside component in a unique corner of its attachment vertex.
Extra loops at a block vertex lie in such corners as well. Group all
parts in each corner into one insert.

If the root is a loop, take that loop as the root block. Its two sides
are its two corners, and the parts on either side form the two inserts.
In both cases a nonempty insert is rooted at the first outgoing dart
in its corner sector, in the oriented cyclic order. Conversely, glue
an arbitrary rooted insert at each corner, identifying its root origin
with the corner vertex and starting its cyclic order with its root dart.
This reconstructs the map and its rotations uniquely. Each insertion
meets the root block at only one vertex, so it cannot enlarge that
block. These constructions are inverse.

Rooting also identifies the corners individually. Indeed, an
automorphism fixing a root dart and preserving vertex rotations and
edge reversal fixes all darts by connectedness. We may thus choose a
deterministic order of the $2m$ corners of every rooted block, with no
symmetry divisor.

For the weight claim, join $M_1$ and $M_2$ at one vertex and write an
edge subset as $A_1\cup A_2$. The number of components and the number
of vertices are, respectively,
\[
 k_{M_1}(A_1)+k_{M_2}(A_2)-1,\qquad
 |V(M_1)|+|V(M_2)|-1.
\]
The exponent $2k+|A|-|V|-1$ is therefore the sum of the two exponents.
Summing independently over $A_1,A_2$ proves multiplicativity.
A size-$m$ root block contributes $b_mz^mT(z)^{2m}$; adding the
trivial map gives \eqref{block:equation}.
\end{proof}

\subsection{An endpoint criterion for finite variance}

An $m$-edge block supplies $2m$ child slots, so the first two offspring
moments will be controlled by the first two derivatives of the block
series. The next lemma extracts the required endpoint identities and
bounds from the word coefficient estimate.

\begin{lemma}[Quadratic inverse criterion]\label{block:critical}
Suppose $T$ and $\Phi$ have nonnegative coefficients and constant
term $1$, and satisfy \eqref{block:equation}. Let $u>0$ satisfy
$T(u)<\infty$, and suppose that, for some $c>0$,
\begin{equation}\label{block:lower-assumption}
 \sum_{n\ge1}na_nu^nx^n\ge c(1-x)^{-1/2}
 \qquad\text{for all $x<1$ sufficiently close to $1$}.
\end{equation}
Put $\tau=T(u)$ and $v=u\tau^2$. Then
\begin{equation}\label{block:critical-equations}
 \Phi(v)=\tau,\qquad
 2v\Phi'(v)=\Phi(v),\qquad
 \Phi''(v)<\infty,
\end{equation}
where endpoint derivatives denote the corresponding nonnegative series.
\end{lemma}

\begin{proof}
Integrating \eqref{block:lower-assumption} gives, for $z<u$ close to $u$,
\begin{equation}\label{block:square-root-gap}
 \tau-T(z)
 =\int_{z/u}^{1}\frac1x
        \sum_{n\ge1}na_nu^nx^n\,\dd x
 \ge c_1\sqrt{u-z}.
\end{equation}
All integrands are nonnegative, so monotone convergence justifies the
endpoint integration. The map
\[
 y(z)=zT(z)^2
\]
is strictly increasing from $[0,u)$ onto $[0,v)$. Nonnegative series
substitution extends \eqref{block:equation} to these real arguments and
then, by monotone convergence, to $\Phi(v)=\tau$. In particular $\Phi$
is finite at $v$ and analytic below it.

The function $f(y)=y/\Phi(y)^2$ is the inverse of $y(z)$ on these
intervals. Hence
\[
 f'(y)=\frac{\Phi(y)-2y\Phi'(y)}{\Phi(y)^3}\ge0
 \qquad(0<y<v).
\]
It follows that $\Phi'(y)\le\Phi(y)/(2y)$ near $v$. Since the derivative
series has nonnegative coefficients, $\Phi'(v-)$ exists and is finite.
Furthermore,
\[
 v-y(z)=(u-z)\tau^2+z\big(\tau^2-T(z)^2\big)
       \ge c_2\big(\tau-T(z)\big)
\]
near $u$. Together with \eqref{block:square-root-gap}, this yields
\begin{equation}\label{block:quadratic}
 0\le u-f(y)\le C(v-y)^2\qquad(y\uparrow v).
\end{equation}
The displayed formula for $f'$ has a finite limit at $v$; the secant
bound \eqref{block:quadratic} forces that limit to be zero. This proves
$2v\Phi'(v)=\Phi(v)$.

For $y<v$,
\[
 f''(y)
  =-\frac{4\Phi'(y)}{\Phi(y)^3}
   -\frac{2y\Phi''(y)}{\Phi(y)^3}
   +\frac{6y\Phi'(y)^2}{\Phi(y)^4}.
\]
If $\Phi''(v)=\infty$, the middle term would force $f''(y)\to-\infty$.
For every $K>0$, sufficiently near $v$ we would then have
$f''(y)\le-K$. Integrate first on an interior interval $[y,r]$
and let $r\uparrow v$, using $f'(v-)=0$, to obtain
$f'(y)\ge K(v-y)$. A second integration, using $f(v-)=u$, gives
\[
 u-f(y)\ge\frac K2(v-y)^2
\]
there, contradicting \eqref{block:quadratic} when $K>2C$. Thus
$\Phi''(v)<\infty$.
\end{proof}

By Corollary~\ref{word:catalan}, $T(u)<\infty$.
Proposition~\ref{ladder:lower} supplies \eqref{block:lower-assumption},
so Lemma~\ref{block:critical} applies. Notice that it does not require
analytic continuation of $\Phi$ across $v$.

\subsection{The marked branching law}

The exponential tilt below is the standard passage from simply
generated trees to a conditioned Galton--Watson law; see
Janson~\cite[Section~4]{Janson2012}. The preceding lemma supplies the
criticality and moment properties for these particular FK weights.

Define a distribution supported on even integers by
\begin{equation}\label{block:offspring}
 p_{2m}=\frac{b_mv^m}{\Phi(v)}\quad(m\ge0),\qquad
 b_0=1,\qquad p_{2m+1}=0.
\end{equation}
If $\xi$ has this distribution, then
\begin{equation}\label{block:moments}
 \E\xi=1,\qquad
 \sigma^2:=\operatorname{Var}(\xi)
     =1+\frac{4v^2\Phi''(v)}{\Phi(v)}\in(0,\infty).
\end{equation}
Here normalization, criticality, and finiteness follow from
\eqref{block:critical-equations}; positivity also follows from
$p_0,p_2>0$.

Construct a plane Galton--Watson tree with offspring law
\eqref{block:offspring}. Independently at each node with $2m>0$
children, assign a rooted block $B$ of size $m$ with probability
$a(B)/b_m$. The children occupy the block's ordered corners. A leaf
represents a trivial insert. Iterating the corner substitutions recovers
a rooted map. We call the origin of a block's root dart its \emph{pole};
the pole of a trivial insert is its attachment vertex.

\begin{proposition}[Exact finite representation]\label{block:finite-law}
The marked tree conditioned to have $N=2n+1$ nodes produces exactly the
map marginal of \eqref{eq:finite-law}. Every such size is attainable.
\end{proposition}

\begin{proof}
A block of size $m$ gives $2m$ children. Thus a map with $n$ edges
corresponds to a tree with $N-1=2n$ child slots. The probability factor
for a node carrying a prescribed block $B$ is
\[
 \frac{b_mv^m}{\Phi(v)}\frac{a(B)}{b_m}
 =\frac{v^ma(B)}{\Phi(v)}.
\]
For a leaf it is $1/\Phi(v)$. Multiplication over all nodes and
Lemma~\ref{block:decomposition} give the unconditioned probability
\[
 \Phi(v)^{-N}v^na(M).
\]
At fixed $N=2n+1$ the first two factors are constant, so the conditional
map law is proportional to $a(M)$, as required.
Since $p_0,p_2>0$, a full binary tree with $n$ internal nodes and all
internal blocks chosen as one-link maps has positive probability.
This gives an admissible tree for every $n\ge1$.
\end{proof}

Figure~\ref{fig:corners} illustrates why trivial inserts must be retained
as leaves: they record corners even when they add no map edges.

\begin{figure}[tb]
\centering
\begin{tikzpicture}[>=Stealth,font=\small]
 \begin{scope}[xshift=-3.5cm,yshift=0.4cm]
  \coordinate (u) at (0,0);
  \coordinate (v) at (1.45,0);
  \coordinate (w) at (2.9,0);
  \draw[thick,->] (u)--(v);
  \draw[thick] (v)--(w);
  \draw[thick] (u) .. controls (-0.85,1.05) and (0.85,1.05) .. (u);
  \draw[thick] (w) .. controls (2.05,1.05) and (3.75,1.05) .. (w);
  \foreach \x in {u,v,w}{\fill (\x) circle (1.8pt);}
  \node[below=4pt] at (u) {$u$};
  \node[below=4pt] at (v) {$v$};
  \node[below=4pt] at (w) {$w$};
  \node at (1.45,-0.95) {$(\deg(u),\deg(v),\deg(w))=(3,2,3)$};
  \node at (1.45,1.45) {map: four edges};
 \end{scope}
 \begin{scope}[xshift=2.8cm,yshift=1.55cm,
  block/.style={circle,draw,fill=black!10,inner sep=2pt,minimum size=15pt},
  leaf/.style={circle,draw,inner sep=1.5pt,minimum size=13pt}]
  \node[block] (r) at (0,0) {$u$};
  \node[block] (a) at (-1.15,-0.85) {$u$};
  \node[block] (b) at (1.15,-0.85) {$v$};
  \node[leaf] (a1) at (-1.55,-1.65) {$u$};
  \node[leaf] (a2) at (-0.75,-1.65) {$u$};
  \node[leaf] (b1) at (0.7,-1.65) {$v$};
  \node[block] (b2) at (1.6,-1.65) {$w$};
  \node[leaf] (c1) at (1.2,-2.45) {$w$};
  \node[leaf] (c2) at (2.0,-2.45) {$w$};
  \foreach \x/\y in {r/a,r/b,a/a1,a/a2,b/b1,b/b2,b2/c1,b2/c2}
   {\draw (\x)--(\y);}
  \node[above=2pt] at (r.north) {link};
  \node[left=2pt] at (a.west) {loop};
  \node[right=2pt] at (b.east) {link};
  \node[right=2pt] at (b2.east) {loop};
  \node at (0.1,-3.02) {block tree: nine nodes};
 \end{scope}
\end{tikzpicture}
\caption{Corner substitution and pole multiplicities. Shaded nodes carry
one-edge blocks; white nodes are trivial inserts. Node labels give
their poles in the map. The eight nonroot nodes have pole counts
$3,2,3$, exactly the vertex degrees. The arrow marks the root dart.}
\label{fig:corners}
\end{figure}

